\documentclass[a4paper,11pt]{ltjsarticle}
\usepackage[haranoaji]{luatexja-preset}
\usepackage{amsmath,amssymb}
\usepackage[top=12mm,bottom=10mm,left=14mm,right=14mm]{geometry}
\usepackage{array}
\usepackage{tikz}
\usetikzlibrary{calc}
\pagestyle{empty}
\setlength{\parindent}{0pt}
\newlength{\qcol}\setlength{\qcol}{0.47\textwidth}
\newlength{\qgap}
\newlength{\anscolw}\setlength{\anscolw}{0.30\textwidth}
\newlength{\ansh}
\newcommand{\Q}[2]{\parbox[t]{\qcol}{\raggedright (#1)\hspace{0.6em}$\displaystyle #2$}}
\newcommand{\QR}[4]{\Q{#1}{#2}\hfill\Q{#3}{#4}\par\vspace{\qgap}}
\newcommand{\anscell}[1]{\rule[-\ansdrop]{0pt}{\ansh}(#1)}
\newlength{\ansdrop}

\begin{document}
{\large\bfseries 計算問題　16}\hfill 名前(\underline{\hspace{32mm}})\par\vspace{3mm}
■（1）〜（16）の計算をしなさい。（17）〜（21）は方程式を解きなさい。\par\vspace{3mm}
\setlength{\qgap}{5.4mm}
\QR{1}{-6+(-15)}{2}{(-72)\div(+6)}
\QR{3}{(2a-1)+(3a-5)}{4}{\{(6-14)\div(-2)\}+5\times(-3)}
\QR{5}{(-15)\times(-11)}{6}{3(6a-1)-2(a-7)}
\QR{7}{(16a-12)\div4}{8}{7-3\times5}
\QR{9}{-3\times6-72\div(-2^3)}{10}{\dfrac{x+5}{3}\times12}
\QR{11}{3\div9\times(-27)}{12}{5+6^2+(-7)}
\QR{13}{-120a\div8}{14}{(+9)+(-11)+(-13)}
\QR{15}{\dfrac{1}{2}(4x-8)+\dfrac{1}{3}(9x-3)}{16}{9-(-2)\div\left(-\dfrac{2}{5}\right)}
\QR{17}{-2x-7=-6x+5}{18}{\dfrac{1}{2}x+3=\dfrac{1}{3}x+5}
\QR{19}{0.09x-0.08=1}{20}{5(2x-3)=3(2x-4)}
\vspace{1mm}
\Q{21}{5(x+1)-2=3(x-3)}\par\vspace{3mm}
\setlength{\ansh}{9.5mm}\setlength{\ansdrop}{6mm}%
\begin{center}\begin{tabular}{|p{\anscolw}|p{\anscolw}|p{\anscolw}|}\hline
\anscell{1} & \anscell{2} & \anscell{3} \\\hline
\anscell{4} & \anscell{5} & \anscell{6} \\\hline
\anscell{7} & \anscell{8} & \anscell{9} \\\hline
\anscell{10} & \anscell{11} & \anscell{12} \\\hline
\anscell{13} & \anscell{14} & \anscell{15} \\\hline
\anscell{16} & \anscell{17} & \anscell{18} \\\hline
\anscell{19} & \anscell{20} & \anscell{21} \\\hline
\end{tabular}\end{center}

\end{document}
